% !TEX root = ../zavrsni_rad.tex

\chapter{Zaključak}
\label{pog:zakljucak}

\section{Pregled postignutih rezultata}

U ovom radu razvijen je dvoslojni sustav za planiranje i praćenje putanje mobilnog robota u poznatom statičkom prostoru. Globalni sloj temelji se na A* algoritmu koji pretražuje diskretiziranu mrežu rezolucije $0{,}2$ m uz ukupnu inflaciju prepreka $0{,}4$ m, a dobivena putanja se izglađuje uz vizualizacijski profil brzine prilagođen zakrivljenosti. Lokalni sloj koristi NMPC regulator s proširenim modelom monocikla i SQP\_RTI solverom koji na svakom uzorkovnom koraku od $\Delta t = 0{,}1$ s rješava OCP s horizontom $N = 20$ koraka.

Sustav je testiran na sedam scenarija koji pokrivaju spektar od jednostavne ravne putanje do gusto popunjene okoline i namjerno zadanog poremećaja stanja. Svi scenariji završili su uspješnim dosezanjem cilja, s RMSE od $34{,}4$ do $72{,}4$ mm u statičkim scenarijima i prosječnim CPU vremenom od $0{,}089$ do $0{,}118$ ms po koraku.

\section{Analiza ključnih nalaza}

\subsection*{Preciznost praćenja i utjecaj geometrije okoline}

Rezultati pokazuju da preciznost praćenja nije monotono određena brojem prepreka, nego geometrijom putanje koju prepreke nameću. Scenarij s jednom preprekom (RMSE $= 46{,}9$ mm) postiže lošiji rezultat od gustog rasporeda sedam prepreka (RMSE $= 42{,}9$ mm), jer raspoređene prepreke u potonjem slučaju dopuštaju glatku dijagonalnu trajektoriju, dok jedna prepreka u nepovoljnom položaju prisiljava na oštriji zaokret. Scenariji s hodnicima (L-hodnik, U-oblik) s RMSE od $72{,}4$ mm i $55{,}0$ mm postavljaju najveći zahtjev jer oblik hodnika nameće skretanja blizu $90°$.

\subsection*{Automatski oporavak od poremećaja}

Scenarij perturbacije pokazuje prednost MPC arhitekture: oporavak od poremećaja nije posebno kodiran, već je prirodna posljedica zatvorene optimizacijske petlje. Nagli pomak od $\Delta y = 1{,}0$ m kompenziran je unutar $2{,}2$ s ($22$ koraka) bez ikakvih modifikacija sustava. Sustav na svakom koraku prima stvarno stanje i rješava OCP prema naprijed, pa automatski ispravlja grešku bez posebne logike za oporavak.

\subsection*{Računalna provedivost}

Prosječno CPU vrijeme kreće se od $0{,}089$ do $0{,}118$ ms po koraku, što je daleko ispod uzornog intervala od $100$ ms. Zanimljiv nalaz je da CPU prosjek ne ovisi izravno o broju prepreka: scenarij poremećaja bez prepreka postiže najviši prosjek ($0{,}118$ ms) zbog intenzivnih korektivnih iteracija, dok gusti raspored sedam prepreka iznosi $0{,}107$ ms. Topologija putanje i dinamika stanja imaju veći utjecaj na računalni teret od broja mekih ograničenja.

\subsection*{Analiza osjetljivosti parametara}

Analiza horizonta $N$ otkriva zasićenje preciznosti za $N \geq 10$ (RMSE $\approx 43$--$47$ mm), dok CPU raste gotovo linearno. Odabrani $N = 20$ ($0{,}305$ ms/korak) postiže $33$\% manji CPU od $N = 30$ uz manje od $3$ mm razlike u RMSE. Analiza težina $Q/R$ pokazuje da što je $Q$ veći u odnosu na $R$, robot preciznije prati putanju. Konfiguracija Q$<$R rezultira najlošijim RMSE ($272{,}9$ mm) jer robot tada više kažnjava agresivno skretanje nego grešku pozicije.

\section{Ograničenja implementiranog sustava}

Implementirani sustav počiva na nekoliko pretpostavki koje ograničavaju izravnu primjenjivost u stvarnim scenarijima:

\textbf{Poznata i statička okolina.} A* algoritam zahtijeva potpunu kartu okoline unaprijed. U dinamičkim okolinama s pokretnim preprekama potrebna je integracija senzorskog sustava i reaktivno replaniranje.

\textbf{Kinematički model bez dinamike.} Model monocikla zanemaruje inerciju i trenje kotača. Na visokim brzinama ili klizavim podlogama ova aproksimacija nije valjana i zahtijeva dinamički model.

\textbf{Kružne prepreke.} Implementirana ograničenja prepreka pretpostavljaju kružni oblik. Prepreke proizvoljnog oblika zahtijevaju konveksnu aproksimaciju poligona ili pristup temeljen na polju predznačenih udaljenosti (SDF).

\textbf{Simulacijsko okruženje i identičan model.} Svi rezultati dobiveni su u simulaciji. Dodatno, NMPC predikcijski model i model koji simulira stvarnog robota su identični --- oba koriste isti kinematički model monocikla. Zbog toga NMPC "zna" točno kako će se robot ponašati, što daje gornju granicu preciznosti: izmjerene greške praćenja predstavljaju optimistični scenarij koji se na stvarnom robotu ne može dosegnuti. Na stvarnom hardveru prisutni su šum senzora, kašnjenje aktuatora i trenje koje uzrokuju odstupanje od predikcijskog modela, pa je stvarna preciznost praćenja nešto lošija.

\section{Smjernice za buduća istraživanja}

Na temelju postignutih rezultata i identificiranih ograničenja, predlažu se sljedeće smjernice za daljnji razvoj:

\begin{enumerate}
  \item \textbf{Dinamičke prepreke i reaktivno replaniranje:} implementacija A* ponovnog planiranja putanje u stvarnom vremenu pri detekciji novih prepreka
  \item \textbf{Dinamički model robota:} uključivanje inercije i trenja u model kako bi simulacija bila bliža stvarnom ponašanju robota
  \item \textbf{Contouring MPC:} varijanta NMPC-a koja kažnjava bočno odstupanje i brzinu napredovanja po putanji umjesto praćenja diskretnih referentnih točaka
  \item \textbf{Adaptivno podešavanje:} automatsko prilagođavanje težina $Q$ i $R$ ovisno o trenutnoj situaciji, npr.\ agresivnije u zavojima i glađe na ravnom dijelu putanje
  \item \textbf{Testiranje na stvarnom hardveru:} pokretanje sustava na Raspberry Pi ili sličnom računalu kako bi se potvrdilo da radi dovoljno brzo u stvarnim uvjetima; pri tome bi trebalo implementirati i transformaciju $(v, \omega) \to (v_L, v_R)$ opisanu u odjeljku~\ref{sec:diff_drive} za naredbu motornim kontrolerima diferencijalnog robota
\end{enumerate}

\section{Završna ocjena}

Kombinacija A* algoritma i NMPC-a pokazala se kao dobro rješenje za navigaciju mobilnog robota --- sustav precizno prati putanju, automatski se oporavlja od poremećaja i radi unutar zahtjeva intervala uzorkovanja od 100 ms s rezervom od dva reda veličine. Za razliku od jednostavnijih regulatora koji reagiraju tek nakon što se greška dogodi, NMPC gleda unaprijed i pokušava spriječiti grešku dok se još može.

Alat acados znatno pojednostavljuje implementaciju jer automatski generira efikasan solver iz zadane OCP formulacije, što oslobađa korisnika ručnog izvođenja derivacija i pisanja vlastitog solvera. 